% Generated draft: requires mathematical and source review.
\documentclass[11pt]{article}
\usepackage[margin=25mm]{geometry}
\usepackage{fontspec}
\setmainfont{Latin Modern Roman}
\usepackage{amsmath,amssymb,mathtools,unicode-math}
\setmathfont{Latin Modern Math}
\usepackage{enumitem,needspace,xurl,hyperref}
\setlength{\parindent}{0pt}
\setlength{\parskip}{5pt}
\setlength{\emergencystretch}{3em}
\newcommand{\N}{\mathbb{N}}
\newcommand{\Z}{\mathbb{Z}}
\newcommand{\Q}{\mathbb{Q}}
\newcommand{\R}{\mathbb{R}}
\newcommand{\C}{\mathbb{C}}
\newcommand{\F}{\mathbb{F}}
\newcommand{\A}{\mathbb{A}}
\newcommand{\PP}{\mathbb{P}}
\newcommand{\E}{\mathbb{E}}
\newcommand{\eps}{\varepsilon}
\newcommand{\dd}{\,\mathrm{d}}
\begin{document}
\section*{UnsolvedMath ID: 30001307}
\textbf{Source problem code:} catalogue entry
\subsection{Definitions and mathematical statement}
Let \(\Omega\subset\R^3\) be a bounded multiply connected domain. The supplied record does not specify its boundary regularity or the number of its boundary components. Whenever the boundary regularity and boundary data permit the following integrals, write \(\Gamma_j\), \(j\in J\), for its boundary components, let \(n\) be the outward unit normal, and define the flux of a prescribed boundary velocity \(a\) through \(\Gamma_j\) by
\[
\Phi_j(a):=\int_{\Gamma_j}a\cdot n\,\dd S.
\]
The condition explicitly stated in both supplied formulations is
\[
\Phi_j(a)=0\qquad\text{for every }j\in J.
\]
This requires zero flux separately through every boundary component; it does not allow arbitrary nonzero component fluxes that merely cancel in their sum. When there are finitely many boundary components, the latter, distinct condition is
\[
\sum_{j\in J}\Phi_j(a)=0.
\]
The record calls the componentwise condition the ``general flux condition,'' but does not explain how it differs from the ``restricted flux condition'' in the referenced existence theorem. Consequently, neither the intended meaning of ``arbitrary flux'' nor the intended extension is determined unambiguously.

Under the conventional prescribed-boundary-velocity interpretation, the stationary incompressible Navier--Stokes system would be written
\[
\begin{cases}
-\nu\Delta u+(u\cdot\nabla)u+\nabla p=f & \text{in }\Omega,\\
\nabla\cdot u=0 & \text{in }\Omega,\\
u=a & \text{on }\partial\Omega,
\end{cases}
\]
where \(\nu>0\) is the viscosity, \(u\) is the velocity, \(p\) is the pressure, and \(f\) denotes a body force if one is included in the intended problem. This display supplies notation for a possible interpretation, not hypotheses absent from the record: the source does not specify the boundary condition explicitly, whether a force is present, the data spaces, or the required solution class.

\textbf{Source question, with unresolved qualifications.} Does the cited existence theorem for the stationary three-dimensional Navier--Stokes equations on bounded multiply connected domains extend from its unspecified restricted flux hypothesis to every admissible choice of data satisfying the intended general flux hypothesis? Read literally, the supplied target flux hypothesis is \(\Phi_j(a)=0\) for every \(j\). An alternative target allowing nonzero component fluxes subject only to zero total boundary flux is not asserted here. A fully quantified mathematical conjecture requires identification of the restricted theorem, resolution of this flux distinction, and specification of the admissible domains, data, and solution class.
\subsection{Short English statement}
The supplied question asks whether an existence theorem for stationary three-dimensional Navier--Stokes flow extends beyond a restricted flux hypothesis. Its wording requires zero flux through each boundary component, while the intended distinction between restricted and general flux is unspecified. Boundary conditions, regularity assumptions, forcing, and solution class also require clarification before a precise universal existence claim can be stated.
\subsection{Sources}
\begin{itemize}
\item Supplied local UnsolvedMath registry record, permanent ID \(30001307\), exact source problem code \texttt{catalogue entry}, title: ``General-Flux Existence for Stationary Three-Dimensional Navier--Stokes.'' Collection: Oberwolfach Reports Open Problems. No individual proposer is identified; the recorded proposed year is 2009.
\item Supplied report citation: \emph{Mathematical Aspects of Hydrodynamics} (2010), Oberwolfach Reports, pp.\ 1921--1942, DOI 10.4171/owr/2009/34. \url{https://doi.org/10.4171/owr/2009/34}. The record describes these reports as documenting mathematical workshops and their open-problem discussions. The supplied report year and DOI are retained as recorded.
\item Supplied literature reference labeled ``Kozono--Yanagisawa theorem'': \url{https://doi.org/10.1007/s00209-008-0361-2}.
\item Supplied literature reference labeled ``Farwig--Kozono--Yanagisawa theorem'': \url{https://doi.org/10.1007/s00208-011-0716-6}.
\item Literature assessment dated exactly \texttt{2026-08-21}: the supplied curation labels the problem ``partially solved'' and states that existence under the general flux condition has been proved under smallness or geometric hypotheses, while the unrestricted three-dimensional Leray problem is not completely settled. This is supplied dated curation derived from the OWR record, not an independently verified literature assessment. The cited claims and time-sensitive status require independent verification.
\item The registry labels the statement \texttt{corrected\_verified} and describes its verification as meaning-preserving semantic, typographic, and LaTeX cleanup without unsupported source-dependent additions. These are supplied metadata, not independent verification of the mathematical formulation or cited results.
\end{itemize}
\end{document}
